\section{OP--DT primary-source comparison}
\label{sec:op-dt-primary-source-comparison}

\subsection{Claim attacked}

The claim under attack is the unqualified identification
\begin{equation}
\label{eq:opdt-claim-attacked}
Z^{X}_{\mathrm{HilbDT}}=Z^{X}_{\mathrm{OP}}.
\end{equation}
If \eqref{eq:opdt-claim-attacked} holds in the normalisation of this
paper, then
\begin{equation}
\label{eq:opdt-cdt-value-if}
C_{\mathrm{DT}}=C_{\mathrm{OP}}=-4096.
\end{equation}
The primary sources below prove \eqref{eq:opdt-cdt-value-if} only after
replacing \(Z^{X}_{\mathrm{HilbDT}}\) by the Oberdieck--Pixton reduced
Gromov--Witten, equivalently reduced stable-pairs, branch.

\subsection{Oberdieck--Pixton branch}

Let \(X=S\times E\), and let
\begin{equation}
\label{eq:opdt-class}
(\beta_h,d)=\iota_{S*}\beta_h+\iota_{E*}(d[E]),
\qquad
\langle\beta_h,\beta_h\rangle=2h-2,
\end{equation}
with \(\beta_h\) primitive.  Oberdieck--Pixton use the reduced
disconnected Gromov--Witten invariant
\begin{equation}
\label{eq:opdt-op-invariant}
\mathsf N^X_{g,h,d}
=
\int_{[\overline M^{\bullet}_{g,1}(X,(\beta_h,d))]^{\mathrm{red}}}
\operatorname{ev}_1^*
\left(\pi_1^*\beta_h^\vee\cup\pi_2^*[\mathrm{pt}]\right),
\qquad
\langle\beta_h,\beta_h^\vee\rangle=1.
\end{equation}
Their primitive partition function is
\begin{equation}
\label{eq:opdt-op-series}
\mathcal Z_{\mathrm{OP}}(u,Q,T)
=
\sum_{g\geq0}\sum_{h\geq0}\sum_{d\geq0}
\mathsf N^X_{g,h,d}\,u^{2g-2}Q^{h-1}T^{d-1}.
\end{equation}
Here
\begin{equation}
\label{eq:opdt-op-variables}
P=e^{iu}=e^{2\pi iz},\qquad
Q=e^{2\pi i\tau},\qquad
T=e^{2\pi i\widetilde\tau}.
\end{equation}
Theorem~1 of Oberdieck--Pixton \cite{OB} states
\begin{equation}
\label{eq:opdt-op-theorem}
\mathcal Z_{\mathrm{OP}}(u,Q,T)
=
-\frac{1}{\chi_{10}^{\mathrm{OP}}(P,Q,T)}.
\end{equation}

The product in Oberdieck--Pixton \cite[(3)]{OB} is
\begin{equation}
\label{eq:opdt-op-chi}
\chi_{10}^{\mathrm{OP}}(P,Q,T)
=
PQT
\prod_{k,h,d}
\left(1-P^kQ^hT^d\right)^{c_{\mathrm{K3}}(4hd-k^2)},
\end{equation}
where
\begin{equation}
\label{eq:opdt-op-product-domain}
k\in\mathbb Z,\qquad h,d\geq0,\qquad
(h>0\ {\rm or}\ d>0)\ {\rm or}\ (h=d=0,\ k<0).
\end{equation}
The coefficient function is fixed by the K3 elliptic genus
\begin{equation}
\label{eq:opdt-k3-coefficients}
Z_{\mathrm{K3}}(z,\tau)
=
\sum_{h\geq0,\ k\in\mathbb Z}
c_{\mathrm{K3}}(4h-k^2)P^kQ^h
=2\phi_{0,1}(\tau,z).
\end{equation}
If
\begin{equation}
\label{eq:opdt-phi-coefficients}
\phi_{0,1}(\tau,z)=\sum_{n\geq0,\ l\in\mathbb Z}f(n,l)q^nr^l,
\end{equation}
then
\begin{equation}
\label{eq:opdt-coeff-doubling}
c_{\mathrm{K3}}(4hd-k^2)=2f(hd,k).
\end{equation}
The monic Borcherds product in this paper is
\begin{equation}
\label{eq:opdt-d5-product}
D_5(Z)
=
q^{1/2}r^{1/2}s^{1/2}
\prod_{(n,l,m)\in\Gamma_{\mathrm{eff}}}
\left(1-q^nr^ls^m\right)^{f(nm,l)}
=64^{-1}\Delta_5(Z),
\end{equation}
with
\begin{equation}
\label{eq:opdt-borch-vars}
q=Q,\qquad r=P,\qquad s=T.
\end{equation}
Equations \eqref{eq:opdt-op-chi}--\eqref{eq:opdt-borch-vars} give
\begin{equation}
\label{eq:opdt-chi-delta}
\chi_{10}^{\mathrm{OP}}(P,Q,T)
=D_5(Z)^2
=4096^{-1}\Delta_5(Z)^2.
\end{equation}
Therefore the OP branch is
\begin{equation}
\label{eq:opdt-op-constant}
\mathcal Z_{\mathrm{OP}}(u,Q,T)
=-4096\,\Delta_5(Z)^{-2},
\qquad
C_{\mathrm{OP}}=-4096.
\end{equation}

\subsection{Reduced stable-pairs bridge already proved for OP}

Oberdieck--Pandharipande use stable-pairs variables
\begin{equation}
\label{eq:opdt-opand-pair-vars}
y=-P=-e^{iu}.
\end{equation}
For fixed \(n,\beta,d\) they define reduced stable-pairs invariants
\begin{equation}
\label{eq:opdt-pair-invariant}
\mathsf P^X_{n,\beta,d}
=
\int_{[P_n(X,(\beta,d))]^{\mathrm{red}}}
\tau_0\!\left(\pi_1^*\beta^\vee\cup\pi_2^*([0])\right)
\end{equation}
and the series
\begin{equation}
\label{eq:opdt-pair-series-fixed}
\mathsf P^X_{\beta,d}(y)
=
\sum_{n\in\mathbb Z}\mathsf P^X_{n,\beta,d}y^n.
\end{equation}
For primitive \(\beta_h\), Oberdieck--Pandharipande prove the
Gromov--Witten/pairs comparison \cite[Proposition after Conjecture D]{OPand}
\begin{equation}
\label{eq:opdt-gw-pair}
\sum_{n\in\mathbb Z}\mathsf P^X_{n,\beta_h,d}y^n
=
\sum_g\mathsf N^{X\bullet}_{g,\beta_h,d}u^{2g-2},
\qquad
y=-e^{iu}.
\end{equation}
Oberdieck--Pixton use the same reduced stable-pairs branch
\cite[\S4.7]{OB}; in their notation
\begin{equation}
\label{eq:opdt-op-pt}
\mathsf P_{n,h,d}
=
\int_{P_n(X,(\beta_h,d))/E}\nu\,de,
\qquad
\sum_{n\in\mathbb Z}\mathsf P_{n,h,d}y^n
=
\sum_g\mathsf N_{g,h,d}u^{2g-2}.
\end{equation}
Thus the primary-source theorem is
\begin{equation}
\label{eq:opdt-pt-equals-op}
Z^X_{\mathrm{PT,red}}(P,Q,T)
:=
\sum_{h,d,n}\mathsf P_{n,h,d}\,Q^{h-1}T^{d-1}(-P)^n
=
\mathcal Z_{\mathrm{OP}}(u,Q,T).
\end{equation}
Combining \eqref{eq:opdt-op-constant} and
\eqref{eq:opdt-pt-equals-op},
\begin{equation}
\label{eq:opdt-pt-constant}
Z^X_{\mathrm{PT,red}}(P,Q,T)
=-4096\,\Delta_5(Z)^{-2}.
\end{equation}

\subsection{Bryan Hilbert-scheme DT branch}

Bryan defines ordinary ideal-sheaf DT invariants by
\begin{equation}
\label{eq:opdt-bryan-ordinary-dt}
DT_{\beta,n}(X)
=
\int_{\operatorname{Hilb}^{\beta,n}(X)}\nu\,de.
\end{equation}
For \(X=S\times E\) these invariants vanish because the \(E\)-translation
action is fixed-point free.  Bryan therefore defines reduced quotient
Hilbert-scheme invariants \cite[Definition 2.1]{BRYAN}
\begin{equation}
\label{eq:opdt-bryan-reduced-dt}
\mathsf{DT}^{\mathrm{Hilb}}_{\beta+dE,n}(X)
=
\int_{\operatorname{Hilb}^{\beta+dE,n}(X)/E}\nu\,de.
\end{equation}
For primitive \(\beta_h\), write
\begin{equation}
\label{eq:opdt-bryan-index}
\mathsf{DT}^{\mathrm{Hilb}}_{h,d,n}(X)
:=
\mathsf{DT}^{\mathrm{Hilb}}_{\beta_h+dE,n}(X).
\end{equation}
Bryan's partition function is
\begin{equation}
\label{eq:opdt-bryan-series}
\mathsf{DT}^{\mathrm{Bryan}}(X)
=
\sum_{h\geq0}\sum_{d\geq0}\sum_{n\in\mathbb Z}
\mathsf{DT}^{\mathrm{Hilb}}_{h,d,n}(X)\,
\widetilde q_{\mathrm B}^{\,h-1}q_{\mathrm B}^{d-1}(-p_{\mathrm B})^n.
\end{equation}
Bryan states explicitly that his \(q_{\mathrm B},\widetilde q_{\mathrm B}\)
are opposite to Oberdieck--Pandharipande's variables.  Hence
\begin{equation}
\label{eq:opdt-bryan-variable-map}
p_{\mathrm B}=P,\qquad
\widetilde q_{\mathrm B}=Q,\qquad
q_{\mathrm B}=T.
\end{equation}
Equivalently, for this paper's displayed DT variables,
\begin{equation}
\label{eq:opdt-manuscript-dt-map}
q_{\mathrm{DT}}=T=s,\qquad
t_{\mathrm{DT}}=Q=q,\qquad
p_{\mathrm{DT}}=P=r.
\end{equation}

Bryan's Donaldson--Thomas Igusa statement is Conjecture~2.3
\cite{BRYAN}:
\begin{equation}
\label{eq:opdt-bryan-conjecture}
\mathsf{DT}^{\mathrm{Bryan}}(X)
=-\frac1{\chi_{10}(p_{\mathrm B},q_{\mathrm B},
\widetilde q_{\mathrm B})}.
\end{equation}
Under \eqref{eq:opdt-bryan-variable-map} and the symmetry
\(\chi_{10}(P,T,Q)=\chi_{10}(P,Q,T)\),
\eqref{eq:opdt-bryan-conjecture} is exactly
\begin{equation}
\label{eq:opdt-bryan-conj-op-form}
\mathsf{DT}^{\mathrm{Bryan}}(X)
=
\mathcal Z_{\mathrm{OP}}(u,Q,T)
=
-4096\,\Delta_5(Z)^{-2}.
\end{equation}
Bryan proves only the \(h=0\) and \(h=1\) Jacobi coefficients, and the
Behrend-weighted computation assumes the Behrend-function conjecture of
Bryan--Kool:
\begin{equation}
\label{eq:opdt-bryan-h01}
\mathsf{DT}^{\mathrm{Bryan}}_{0}(X)=\frac1{F^2\Delta},
\qquad
\mathsf{DT}^{\mathrm{Bryan}}_{1}(X)=-24\,\frac{\wp}{\Delta}.
\end{equation}
These are the \(\widetilde q_{\mathrm B}^{-1}\) and
\(\widetilde q_{\mathrm B}^{0}\) coefficients of
\(-\chi_{10}^{-1}\), not the full equality
\eqref{eq:opdt-bryan-conj-op-form}.

\subsection{Exact theorem needed}

The missing theorem is the reduced ideal-sheaf/stable-pairs comparison
for \(S\times E\), in the quotient-by-\(E\) Behrend-weighted
normalisation:
\begin{equation}
\label{eq:opdt-needed-coeff}
\mathsf{DT}^{\mathrm{Hilb}}_{h,d,n}(X)
=
\mathsf P_{n,h,d}
\qquad
\hbox{for all }h,d\geq0,\ n\in\mathbb Z,
\end{equation}
with the variable convention
\begin{equation}
\label{eq:opdt-needed-vars}
y=-P,\qquad
(p_{\mathrm B},q_{\mathrm B},\widetilde q_{\mathrm B})
=(P,T,Q).
\end{equation}
Equivalently, the needed theorem is the partition identity
\begin{equation}
\label{eq:opdt-needed-series}
\sum_{h,d,n}
\mathsf{DT}^{\mathrm{Hilb}}_{h,d,n}(X)
Q^{h-1}T^{d-1}(-P)^n
=
\sum_{h,d,n}
\mathsf P_{n,h,d}Q^{h-1}T^{d-1}(-P)^n.
\end{equation}
No non-trivial degree-zero factor is allowed in
\eqref{eq:opdt-needed-series}; if a DT/PT correspondence produces such
a factor \(Z_0(P)\), the manuscript needs the corrected statement
\begin{equation}
\label{eq:opdt-needed-with-factor}
Z^{X}_{\mathrm{HilbDT}}
=
Z_0(P)\,Z^X_{\mathrm{PT,red}},
\qquad
Z_0(P)=1
\quad\hbox{in the reduced quotient normalisation}.
\end{equation}
Equations \eqref{eq:opdt-needed-coeff}--\eqref{eq:opdt-needed-with-factor}
are not proved by \cite{OB}, \cite{OPand}, or \cite{BRYAN}.

\subsection{Strongest true replacement}

The unconditional scalar-square theorem is
\begin{equation}
\label{eq:opdt-true-replacement-op}
Z^X_{\mathrm{OP}}
=Z^X_{\mathrm{PT,red}}
=-4096\,\Delta_5^{-2}.
\end{equation}
The Hilbert-scheme DT scalar square is conditional:
\begin{equation}
\label{eq:opdt-true-replacement-dt}
Z^X_{\mathrm{HilbDT}}
=C_{\mathrm{DT}}\Delta_5^{-2}
\end{equation}
is available only as an assumption, or as a theorem after
\eqref{eq:opdt-needed-series} is supplied.  Under
\eqref{eq:opdt-needed-series},
\begin{equation}
\label{eq:opdt-cdt-under-needed-theorem}
C_{\mathrm{DT}}=-4096.
\end{equation}
Without \eqref{eq:opdt-needed-series}, the source-true statement is
\begin{equation}
\label{eq:opdt-final-status}
C_{\mathrm{OP}}=-4096,
\qquad
C_{\mathrm{DT}}\in\mathbb C^\times\ {\rm separate}.
\end{equation}

\subsection{Audit table}

\[
\begin{array}{c|c|c}
\hbox{object} & \hbox{equation} & \hbox{status} \\
\hline
\hbox{OP reduced GW}
& \eqref{eq:opdt-op-series}
& \hbox{definition in }\cite{OB} \\
\hbox{OP theorem}
& \eqref{eq:opdt-op-theorem}
& \hbox{proved in }\cite[Theorem 1]{OB} \\
\hbox{OP product}
& \eqref{eq:opdt-op-chi}
& \hbox{definition in }\cite[(3)]{OB} \\
\hbox{normalisation conversion}
& \eqref{eq:opdt-chi-delta}
& \hbox{uses }Z_{\mathrm{K3}}=2\phi_{0,1},
\ D_5=64^{-1}\Delta_5 \\
\hbox{OP constant}
& \eqref{eq:opdt-op-constant}
& \hbox{proved: }C_{\mathrm{OP}}=-4096 \\
\hbox{reduced PT branch}
& \eqref{eq:opdt-pt-equals-op}
& \hbox{proved for primitive }\beta_h \\
\hbox{Bryan reduced HilbDT}
& \eqref{eq:opdt-bryan-series}
& \hbox{definition} \\
\hbox{Bryan Igusa formula}
& \eqref{eq:opdt-bryan-conjecture}
& \hbox{conjecture; }h=0,1\hbox{ computed conditionally} \\
\hbox{DT=OP}
& \eqref{eq:opdt-needed-series}
& \hbox{missing theorem} \\
\hbox{manuscript-safe branch}
& \eqref{eq:opdt-final-status}
& \hbox{unconditional}
\end{array}
\]
